% !TeX root = ../papers/nc-p3.tex
\section{滑らかな三交換子の全復元と楕円曲線の負二倍写像}
\label{sec:marked-elliptic-commutator-doubling}
生成添字の三空間を $D$、行列の三空間を $U$ とし、体積形式で
$\bigwedge^2D=D^*$ と読む。$Q:D\to\mathfrak{sl}(U)$ に対し
$C_Q(s\wedge t)=[Q_s,Q_t]$ を考える。滑らかな平面三次曲線 $E\subset\mathbf P(D^*)$
の線型決定式表示 $A(v)$ に対し、その右核点の写像から引き戻した
$\mathcal O(1)$ を $T_A$、平面の線束を $H$、差を $d_A=T_A-H\in\mathrm{Pic}^0(E)$ と書く。
線束の群則は以下では加法で書く。
\begin{theorem}\label{thm:marked-elliptic-commutator-doubling}
$\det C_Q$ が $E$ を定める全ての $Q$ は以下の方法で復元できる。
\begin{enumerate}
\item $d\in\mathrm{Pic}^0(E)\setminus\mathrm{Pic}^0(E)[2]$ を取り、右核線束が
$H+d$ の決定式表示 $A(v)$ を取る。行の線型包の消去空間を $W\subset D^*\otimes U$ とする。
次の写像が $W$ に値を持つ条件は $Q$ に線型で、解空間は六次元である。
\[
 \Gamma_Q:\mathrm{Sym}^2 U^*\longrightarrow D^*\otimes U,\qquad
          B\longmapsto\big(\mathrm{ax}(BQ_i-Q_i^t B)\big)_{i=1}^3 .
\]
ここで $\mathrm{ax}$ は交代形式を体積形式でベクトルにする写像である。
$\Gamma_Q$ が階数6となる非空な開部分の全てが解である。
逆に $d$ が非零二等分点ならこの開部分は空である。
\item 得られた三交換子自身の右核線束は $H-2d$ である。
従って指定した滑らかな三行三列の決定式表示 $C$ を\emph{独立な左右の基底まで}実現する
補助表示の類は、$-2d=d_C$ の相異なる四解ちょうどである。
対称三平面の部分空間を既に固定した場合には、この四選択・六次元線型解の開条件と
最後の基底の線型比較で、順序付き基底と三交換子解を全て得る。
\end{enumerate}
これは印付き二切断方程式の記述であり、各解の四辺有限性および全次数AS性にはなお判定が要る。
\end{theorem}
\begin{proof}[三交換子から補助表示へ]
$\Gamma_Q$ の核に非零 $B$ があれば、$B$ が特異なら $Q_i$ は共通の直線または平面を保ち、
全ての交換子は特異である。$B$ が可逆でも交換子は奇数階の交代形式に対応して特異である。
従って所要の $Q$ では $W=\operatorname{im}\Gamma_Q$ は六次元である。
その消去三平面を三行の双一次形式 $A$ と書く。$v=s\wedge t$ の点で $A(v)u=0$、
$u\ne0$ という条件は、$v\otimes u\in W$、すなわち $Q_s,Q_t$ の両方を
自己随伴にする非零対称形式の存在である。
定理\ref{thm:rank-eight-finite-symmetry}の二行列の対称化行列式により、
これは $\det[Q_s,Q_t]=0$ と同値である。従って二つの三次式は非零定数倍である。
逆向きにも、指定した $A$ の $W$ と $\Gamma_Q$ の像が一致すれば同じ議論が使える。

$\Gamma_Q$ が $W$ に入る線型条件の解次元を示す。変数 $x\in U^*$ について
$Z(x)=\Gamma_Q(xx^t)$ と書く。その各行 $Z_i(x)$ は $x$ との対合が零となる
三成分二次式である。逆にそのような全ての行は一意な跡零 $Q_i$ から来る。
実際 $x\times Q_i^t x$ の符号倍であり、恒等零なら $Q_i$ はスカラー、
二次の三成分から三次式への $x$ との積は全射で次元差は $18-10=8$ である。
従って求める空間は
\[
 W\otimes\mathrm{Sym}^2 U\ \longrightarrow\ D^*\otimes\mathrm{Sym}^3 U
\]
の核である。この積写像は全射である。双対余核の元は、二階偏微分が全て
$W^\perp$ に入る $D$ 値の対称三テンソルである。ある $v_0$ で $W^\perp\to U^*$ の評価が
可逆なので、他の二方向の評価を二つの自己準同型 $B_1,B_2$ で書く。
この対は既約である。共通不変部分空間があれば $\det(v_0'I+v_1'B_1+v_2'B_2)$ が
ブロックで因数分解してしまう。$v_0$ 成分の三テンソルを $P$ とすれば、
他の二成分は $P$ の一因子に $B_i$ を掛けたもので、対称性によりどの因子に掛けても等しい。
$P$ の最小支持空間は両作用で不変となる。さらに二作用を別々の因子に移せるので
$([B_1,B_2]\otimes I\otimes I)P=0$ である。例えば $B_1B_2$ の右の作用を第二因子へ、
左の作用を第三因子へ移してから逆順に第一因子へ戻せる。
非零なら支持が全空間になり二作用が可換となって既約性に反する。
よって余核は零、核次元は $6\cdot6-3\cdot10=6$ である。
\end{proof}

\begin{proof}[決定式表示の座標]
滑らかな決定式では各曲線点で行列の階数は2である。
左核点の線束は $\mathcal O(-1)^3\to\mathcal O^3$ の余核線束 $T'$ であり、
Euler標数より次数3、自由表示を一つ捻れば $H^0(T'-H)=0$ である。
転置により右核点にも同じことが成り立つ。余因子が共通零点のない二次式で
右・左核の積となるので $T_A+T'=2H$ である。
$T\ne H$ なる任意の次数三線束に対し、二つの平面埋込み $(H,T)$ のグラフの
双一次式は三次元をなす。実際 $H$ の基点自由な二切断と $H^0(T)$ の積は、
鉛筆の完全列と $H^0(T-H)=0$ により $H^0(H+T)$ を張る。
この三つの双一次式を行にすれば表示が得られることを次の座標で確かめる。
特に右核線束が等しい二表示は右核の座標を基底変更した後で同じグラフの双一次式を張り、
独立な左右基底まで一致する。

$E[3]$ の二生成平行移動は $H$ の切断へ射影的に作用する。
その二つの持上げの交換倍率は非自明な三乗根である。可換なら共通固有切断の
次数三因子が自由に作用する九元群で不変となり矛盾する。
立方を1に規格化し、一方を対角、他方を三巡回にした基底を取る。
$T$ は $H$ の平行移動による引戻しなので、同じ二作用の基底が取れる。
$E$ の平面式はこの基底でHesse形になる。対角作用の固有文字は自明でなければならない
（座標点はその不動点で $E$ にない）。三巡回の固有文字も同じくその不動点を用いて自明である。
従って三つの立方の係数は等しく、他は積 $xyz$ だけである。
三関係の空間では対角文字を三巡回が推移的に動かすので各文字が一本であり、
三行を同時巡回させて以下のMoore表示となる。
\[
 A_P(x,y,z)=\begin{pmatrix}
 a x&b z&c y\\ c z&a y&b x\\ b y&c x&a z
 \end{pmatrix},\qquad
 \det A_P=(a^3+b^3+c^3)xyz-abc(x^3+y^3+z^3).
\]
$abc=0$ で行列式が恒等零なら、二つの非零係数の三乗和が零で、
行・列の巡回と対角倍率で外積行列にできる。その右核は平面座標の線型同型となり $T=H$ に反する。
行列式が $xyz$ の非零倍でも $E$ は含まれないので、$abc\ne0$、$P=(a:b:c)\in E$ となる。
逆にこのような $P$ は全て滑らかな決定式を与える。
$P\mapsto d(P)=T_{A_P}-H$ は $E$ から $\mathrm{Pic}^0(E)$ への射である。
実際座標の一つが消える九変曲点でも上の外積表示は曲線上常に階数2なので、相対核線束が取れる。
$O=(0:-1:1)$ では核写像は恒等である。この変曲点を原点にすれば、原点を保つ楕円曲線間の射は
群準同型なので $d(-2P)=-2d(P)$ である。
\end{proof}

\begin{proof}[半周期の除外と符号付き二倍]
$d(P)$ が位数2なら左右核線束が一致する。定数基底変更で曲線上の左右核点を
点ごとに揃えると、行列とその転置の差は同じ核を消す。差は交代行列なので、
恒等零でなければその核写像は平面の線型写像である。像が平面を張るから線型同型となり
$T_{A_P}=H$ に反する。よって定数行列を左から掛けて対称にできる。
$L A_P(v)$ の交代成分を消す九線型方程式の行列式は符号を除き
$(a^3-b^3)(a^3-c^3)(b^3-c^3)$ である。

逆に二つの三乗が等しければ、全添字を3で読み、二添字の逆向き巡回で
係数 $(a,b,c)$ を巡回させ $b^3=c^3$ にできる。
非零成分の添字は $i+j+k=0$、係数は $p_{j-k}$（$p_0=a,p_1=b,p_2=c$）である。
二添字を交換すると添字差の符号が変わる。$b/c$ が三乗根なので、一方の添字の
幾何数列状の対角倍率で任意に選んだ二添字について対称にできる。
特に左右核線束は等しく $2d(P)=0$ である。また二変数側を対称にすれば
$W$ は三次交代行列の全空間と、三独立な対称行列 $B_j$ の包の直和になる。
先の $Z(x)$ の対称部分を $\sum h_j(x) B_j$ とする。$Z(x)x=0$ から
$\sum h_j(x)(x^t B_j x)=0$ である。この三つの二次式は基点自由で、
さもなくば $xx^t\in W^\perp$ が階数1となり、その行を最初に選べば決定式に一次因子が出る。
よって三式は正則列で、二次係数の関係は定数の交代三行列によるKoszul関係に尽きる。
従って三係数 $h_j$ は線型従属で $\Gamma_Q$ は $W$ を張れない。

残りの三乗が相異なる場合の非空性と $d_{C_Q}$ を同時に求める。
$Q_0=\operatorname{diag}(d_0,d_1,d_2)$、$Q_1$ の $(i,i+1)$ 成分を $s_i$、
$Q_2$ の $(i,i-1)$ 成分を $r_i$ とし、他の成分は零とする。自由スカラー $h$ について
\[
\begin{split}
 (s_0,s_1,s_2)&=(-3bc(a^3-b^3),3a^2(b^3-c^3),0)
                  +h(a/b,c^2/b^2,c/a),\\
 (m_0,m_1,m_2)&=(c/b,b/a,a/c),\quad r_i=m_i s_i,\\
 k_i&=r_{i+1}-m_i s_{i-1},\quad
 (d_0-d_1,d_1-d_2,d_2-d_0)=(c k_2/b,b k_0/a,a k_1/c).
\end{split}
\]
最後の三成分の和は零で、跡零条件により $d_i$ は決まる。直接代入で
$\operatorname{im}\Gamma_Q\subset W$ である。
$P'=(a(b^3-c^3),b(c^3-a^3),c(a^3-b^3))$ と置くと、$C_Q(v)$ の行 $i$ は
$A_{P'}(v)$ の行 $-i$ に $s_i k_i/P'_{2i}$ を掛けたものである。
各 $k_i$ は $h$ によらず非零、一般 $h$ で $s_i\ne0$ となる。
\path{computations/verify_moore_commutator_doubling.py} が全不定元を残して
六つの対称行列に対する方程式とこの三行の恒等式を検算する。
$P'$ は $P$ の接線と $E$ の第三交点である。曲線式と接線式への代入でこれを確認でき、
変曲点を除く現在の開集合で $P'=-2P$ である。従って行列式は滑らかで
$\Gamma_Q$ は可逆な像を持ち、$d_{C_Q}=d(P')=-2d(P)$ となる。
同じ $A$ の六次元線型空間の階数6の開部分では、$C_Q$ の右核線束も代数的に動く。
有理空間から楕円曲線への有理写像は一定である（一般の直線に制限し完備化する）。
従って得た値は全ての解で同じである。

最後に二つの対称三平面から計算した基底未指定の混合表示を $C$ とする。$d_C\ne0$ なので四つの $-2$ 逆像はいずれも二等分点でない。
各六空間から上の開条件で $Q$ を取ると独立な左右基底変更 $P,R$ で $C_Q=P C R$ にできる。
これは対称三平面の二つの基底の取り直しにほかならず、交換子の跡条件も満たす。
実際の行列を求める際には $C(v_0)$ が可逆な一点で $R$ を消去し、
$C_Q(v)C_Q(v_0)^{-1} P=P C(v)C(v_0)^{-1}$ を二方向について線型に解けばよい。
可逆解は共通スカラーを除き一意で、非スカラーな自己準同型は先の既約性とSchurの補題に反する。
初めから両基底が指定されている場合には同時共役だけを許して $C_Q(v)P=P C(v)$ を用いる。
\end{proof}
